\subsection{Quantitative comparison and sensitivity}

The paper reports a style-ratio range of 0.4--42.4 for 63 events. The present range is 0.03879--463.9. The lower and upper endpoints differ by -0.3612 and 421.5, respectively. These are distribution-endpoint comparisons, not matched-event errors. Matching the number of events does not establish matching event membership or measurements.

The published distance transition near $D/w=60$--$100$ and slope transition near $0.005$ are comparison targets only. They are not inputs to classification, masks, thresholds or geometry. The present independent geometry coverage cannot test either transition. The published individual-result and geometry tables are not imported.

\begin{center}\begin{tabular}{lrr}\toprule Calculation & Minimum $S_{\mathrm R}$ & Maximum $S_{\mathrm R}$\\\midrule

Primary spatial union & 0.03879 & 463.9\\

Sum of image-part counts & 0.03879 & 464\\

Alternative B14 onset & 0.03879 & 463.9\\

Parent in avulsion years & 0 & 30.86\\

Parent in all selected years & 0 & 24.85\\

\bottomrule\end{tabular}\end{center}

The temporal alternatives retain the same annual thresholds and avulsion windows; only the parent accumulation changes. They bracket specific interpretations of the written instruction to use the same time ranges where possible. They are not selected according to agreement with the paper. Every signature-specific result remains available.

The largest primary ratio illustrates the denominator sensitivity. For B10, the two calculated alternatives are B: 124780 avulsion pixels, 269 parent pixels and a ratio of 463.9; W: 520388 avulsion pixels, 25034 parent pixels and a ratio of 20.79. Minimizing avulsion area can choose a much larger ratio when that filter also detects far fewer parent-channel pixels. Both calculations use the same supplied mask geometry. This is an observed numerical consequence of the stated rule, not evidence that the larger ratio correctly measures floodplain disturbance.

The primary range fails to match the published range. Overlap accounting produces much smaller changes than parent-window selection. The experiment therefore establishes reproducibility of the stated downstream computation but does not establish reproduction of the paper's quantitative result. The missing independent geometry additionally prevents testing the proposed downstream transition.

For a limited matched-event check, Figure 5 explicitly prints $S_{\mathrm R}=0.5$ for V15 and $14.9$ for V11. These two numerical annotations were read only after the calculations, without extracting plotted coordinates or image areas. They enter this comparison section only; the numerical analysis never reads them.

For V15, the recomputed ratio is 2.77, a difference of 2.27 (454.0 percent) relative to the printed value.

For V11, the recomputed ratio is 28.66, a difference of 13.76 (92.3 percent) relative to the printed value.

Both calculations rank V11 above V15, but the magnitudes disagree. The brightness-selected V15 fingerprint is visibly sparse compared with the published depiction. The B12 full-event mosaic covers a larger domain than the published central-part view. Its independently chosen detail and annual views do not recover the original validation window; they demonstrate the calculation only. Display extents were refined after comparison, without changing any numerical input or geometry observation. The published images are absent from the generated figures and repository.
